% Folder 57, batch 10 — archivists' pages 181–200.
% Pass: Opus 5.5 (claude-opus-5-5), 2026-10-03 — first pass, unchecked against the pages by a human.
%
% Skipped: 183 and 194 (cover sheets in his hand: their words are the section
% headings below), 196 and 198 (typed pages on dimension and the chain
% condition, carrying no ink of his), 199 (no writing of its own).
\documentclass[11pt,a4paper]{article}
\input{../preamble/grothendieck}

\folder{57}
\batch{10}
\pages{181}{200}
\foldertitle{Picard : tapuscrit annoté (s.d.), lettres (s.d., 1962).}
\dating{1962-[vers 1968]}
\watermark{Édition de démonstration}

\begin{document}

\section{Diviseurs relatifs et schéma de Picard}

\page{181}
\note{la page commence au milieu d'une phrase, dont le début se trouve avant ce lot ; la notation soulignée deux fois ($\mathbf{Pic}$, $\mathbf{Div}$) est rendue en gras, les lettres soulignées une fois ($\underline{E}$, $\underline{L}$, $\underline{M}$, $\underline{H}$) sont gardées soulignées}
d'un préschéma en Picard \ill{}] \struck{\ill{}}. Utilisant ce résultat, on
voit que si on peut trouver un faisceau \add{\uncertain{inversible}}
\uncertain{fondamental} \struck{\ill{} $\underline{\Lambda}$}
$\underline{\Lambda}$ sur $\mathbf{Pic}_{X/S} \times_S X$ \struck{dans ce faisceau}
[par exemple si $X/S$ a une section \ill{}], alors \struck{$\mathbf{Div}_{X/S}$
s'identifie à}
\[
  \mathbf{Div}_{X/S} \simeq
  \mathbf{Div}^{\underline{\Lambda}}_{\mathbf{Pic}_{X/S} \times_S X \,/\, \mathbf{Pic}_{X/S}} .
\]
Cela permet l'\uncertain{étude} d'un \uncertain{schéma}
\uncertain{diviseurs}, indépendamment de celle de $\mathbf{Pic}_{X/S}$, des
\uncertain{foncteurs} \add{de l'\ill{}} $\mathbf{Div}^{\underline{L}}_{X/S}$
et \uncertain{des} \struck{\ill{} représentant} préschémas qui \ill{} les
représentent.

Supposons maintenant que $X$ est propre \add{et \uncertain{plat}} sur $S$, et que
l'on a un \uncertain{isomorphisme} \uncertain{fonctoriel} :
\[
  \underline{H}^0_f(\underline{L} \otimes_{\underline{O}_S} \underline{M})
  \simeq \mathbf{Hom}_{\mathcal{O}}(\underline{E}, \underline{M})
\]
pour $\underline{M}$ \uncertain{quasi-cohérent} sur $S$ [ce qui est vrai si $S$
est local, et $X$ projectif sur $S$] \struck{\ill{}}. Le \ill{} ainsi \ill{}
faisceau \uncertain{cohérent} $\underline{E}$ sur $S$ \add{\uncertain{définissant}}
\struck{fixe} le « \uncertain{foncteur} » $S' \mapsto \underline{H}^0_{f'}(\underline{L}')$.
Supposons \emph{\uncertain{de plus}} que

\page{182}
les fibres de $X/S$ soient \ill{} \struck{\uncertain{intègres}} \add{primaires}.
Alors on voit que l'on a :

\marginal{Prop.}
\[
  \mathbf{Div}^{\underline{L}}_{X/S} \simeq \underline{P}(\underline{E}),
\]
\emph{fibré projectif associé à $\underline{E}$}. Comme ce dernier est aussi
isomorphe à $\mathrm{Proj}($\struck{$\underline{E}$} $\mathrm{Symm}^{\bullet}(\underline{E}))$,
\uncertain{et} \ill{} fibrés \ill{}, on voit aussi :

\textbf{\emph{Corollaire.}} Supposons $X/S$ \emph{primaire}. Alors
\struck{$\mathbf{Div}^{\underline{L}}_{X/S} \to$} les conditions suivantes sont
équivalentes :
\begin{itemize}
  \item[(i)] $\mathbf{Div}^{\underline{L}}_{X/S}$ est plat sur $S$ en $s$ ;
  \item[(ii)] $\mathbf{Div}^{\underline{L}}_{X/S}$ est \emph{simple} sur $S$ en $s$ ;
  \item[(iii)] pour $n$ grand, $\mathrm{Symm}^n(\underline{E})$ est libre sur $S$ en $s$ ;
  \item[(iv)] pour \add{$k$ \uncertain{convenable},} $n$ grand,
    $\mathrm{Symm}^{n}(\underline{E})$ est loc.\ libre sur $S$ \add{en $s$} pour
    $n$ grand ;
  \item[(v)] $\underline{E}$ est \struck{libre} loc.\ libre sur $S$ en $s$ ;
  \item[(vi)] $\underline{L}$ est \uncertain{cohomologiquement} plat rel.\ $S$ ;
  \item[(vii)] $\underline{H}^0_f(\underline{L}) \otimes k(s) \to H^0(X_s, \underline{L}_{(s)})$
    est \emph{surjectif}.
\end{itemize}
\note{devant (vi) et (vii), à gauche, deux mentions brèves : « \ill{} dif. » et
« \ill{} \ill{} » ; le $\otimes k(s)$ de (vii) est récrit sur une rature et
reste douteux}
\marginal{\struck{\ill{}} \ill{} Diagramme \ill{} \ill{} \ill{} fibrés projectifs}
\note{à gauche de la liste, des flèches doubles ($\Uparrow$, $\Downarrow$) et deux
longues flèches courbes relient les conditions, sans qu'on puisse dire avec
certitude lesquelles ; le texte en marge, écrit en biais, renvoie à ce
diagramme d'implications}

\section{Groupe de Picard des cônes et de leurs anneaux locaux au sommet}
\note{titre pris sur la page 183, feuille de garde de sa main, qui porte seulement : « Groupe de Picard des cônes et de leurs anneaux locaux au sommet »}

\page{184}
\emph{Groupe de Picard des cônes épointés}

$X \subset \mathbb{P}^r_k$, cône projetant (affine) $C$, à \struck{l'origine}
sommet $a$, $C_a = \mathrm{Spec}\, \mathcal{O}_{C,a}$.

$C - a$ est un fibré principal sur $X$, de groupe $\mathbb{G}_m$, dont la classe
est celle $\xi \in \mathrm{Pic}(X)$ de $\mathcal{O}_X(1)$. Donc \emph{si $X$ est
régulier}
\[
  \mathrm{Pic}(C - a) \simeq \mathrm{Pic}(X)/\mathbb{Z}.\xi
\]
Considérons le morphisme naturel
\[
  C_a - a \longrightarrow C - a ,
\]
il induit un homomorphisme canonique
\[
  \mathrm{Pic}(C - a) \longrightarrow \mathrm{Pic}(C_a - a)
  \simeq \varinjlim_{U \ni a} \mathrm{Pic}(U - a)
\]
Je dis que \emph{c'est un isomorphisme si $X$ est régulier}. Comme pour
$U = C - Y$ voisinage de $a$, on a
\[
  \mathrm{Pic}(U - a) \simeq \mathrm{Pic}(C - a)/\mathrm{Im}\, \mathbb{Z}^{(I)}
\]
où $I$ est l'ens.\ des comp.\ irréd.\ de $Y$,

\page{185}
l'assertion précédente équivaut à celle-ci : \emph{pour tout diviseur $D$ sur
$C$ tel que} $a \notin \mathrm{supp}\, D$, \emph{on a} $D \sim 0$ \emph{dans} $C - a$.

Pour le prouver, introduisons le cône projetant projectif $\hat{C}$, et
$E = \hat{C} - a$, qui est un fibré vectoriel sur $X$ [savoir celui associé de
façon covariante à $\mathcal{O}_X(1)$]. Alors $D$ provient d'un diviseur $\hat{D}$
sur $\hat{C}$ tel que $a \notin \mathrm{supp}\, \hat{D}$, et par suite le
\struck{\ill{} de $\hat{D}$} diviseur \struck{$E$} sur $E'$ induit par $\hat{D}$
est \emph{à support propre} sur $X$. Nous allons prouver le

\textbf{\emph{Lemme.}} Soit $X$ préschéma régulier, $E$ un fibré vectoriel
\add{de rang 1} sur $X$ ($E = \mathbb{V}(\mathcal{E})$, où $\mathcal{E}$ est un
faisceau inv.\ sur $X$), $D$ un diviseur sur $E$, tel que $\mathrm{supp}\, D$

\page{186}
soit propre sur $X$. \struck{Alors} Soit $V = E - g(X)$, où $g$ est la section
nulle, alors
\[
  D|V \sim 0 \quad \text{dans } V.
\]

\begin{tikzcd}
E \arrow[d, "f"'] & E' \arrow[l, "h'"'] \arrow[d, "f'"'] \\
X & X' \arrow[l, "h"] \arrow[u, bend right=40, "q"']
\end{tikzcd}

On peut supposer que $D$ est un diviseur irréductible, défini par un
sous-préschéma fermé $X'$ de $E$. \struck{Faisons}

Supposons d'abord $X'$ \emph{régulier}, faisons le changement de base
$h \colon X' \to X$, alors $E'$ sur $X'$ est muni d'une section naturelle $q$
telle que $h'q =$ immersion canonique de $X'$ dans $X$
\note{lire sans doute « dans $E$ »}. On a donc \struck{dans $\mathrm{Pic}$}
[(raisonnement \ill{} \uncertain{dans} $\mathrm{Pic}$ \ill{})]
\[
  D = (h'_* q)_* (1_{X'}) = h'_*(q_*.1_{X'})
\]
\marginal{N.B. Cela \ill{} \ill{} \ill{} \ill{} si $h'$ et $q$ sont propres
\uncertain{ou} \ill{}}
mais comme d'ailleurs que
\[
  q_*(1_{X'}) = f'^*(\xi_{E'})
\]
où $\xi_{E'} \in \mathrm{Pic}(X')$ est la classe \ill{} de $E'/X'$. D'ailleurs
$\xi_{E'} = h^*(\xi_E)$, d'où
\[
  q_*(1_{X'}) = f'^* h^*(\xi_E) = h'^* f^*(\xi_E)
\]
et
\[
  D = h'_* h'^*(f^*(\xi_E)) = f^*(\xi_E)\, h'_*(1_{E'})
\]

\page{187}
Or $h$ est un morphisme fini, \add{de \uncertain{dimension relative} (\ill{}) 0}
\struck{\ill{}}, donc $h'_*(1) = d$, d'où
\[
  D = d\, f^*(\xi_E), -
\]
Par suite
\[
  D|V = d\,(f^*(\xi_E)|V) = 0 \qquad \text{cqfd.}
\]
\marginal{Méthode \uncertain{figurant} dans la \ill{} \ill{} \ill{} se
généralise aux fibrés vectoriels quelconques.}

Dans le cas général, supposant cependant le lieu singulier de $X'$ fermé
(\uncertain{donc} \uncertain{contenu} dans un \struck{\ill{}} fermé
\struck{\ill{}} $Y' \subset X'$, $\neq X'$) on considère $Y = f(Y')$, et on
\uncertain{applique} ce qui précède à $X - Y$ et $D|X - Y$, et on gagne \ill{}

On trouve ainsi, dans le cas présent
\[
  \boxed{\mathrm{Pic}(C - a) \simeq \mathrm{Pic}(C_a - a) \simeq \mathrm{Pic}(X)/\mathbb{Z}\xi}
\]
Dans le cas où $C$ est normal \ill{} \struck{\ill{}} $\mathcal{O}_{C,a}$ est
normal) notons par $\mathrm{Pic}(C_a - a)$ \uncertain{aussi} le groupe
$\mathrm{Cl}(\mathcal{O}_{C,a})$ des classes

\page{188}
de diviseurs de $\mathcal{O}_{C,a}$.

Posons
\[
  S = \text{\struck{$\mathrm{Spec}$}}\ C_a = \mathrm{Spec}\, \mathcal{O}_{C,a}, \qquad
  \hat{S} = \mathrm{Spec}\, \widehat{\mathcal{O}_{C,a}}
\]
nous avons un homomorphisme canonique
\[
  \mathrm{Pic}(S - a) \longrightarrow \mathrm{Pic}(\hat{S} - \hat{a})
\]
[qui, dans le cas $S$ normal, s'interprète comme
\[
  \mathrm{Cl}(S) \longrightarrow \mathrm{Cl}(\hat{S})
\]
]
et on peut se demander si c'est un isomorphisme. \struck{Faisons ici} Pour ceci,
nous allons interpréter l'homomorphisme précédent d'une autre manière. Soit
$\tilde{C}$ obtenu à partir de $C$ en « faisant éclater l'origine », donc
$\tilde{C}$ est un fibré vectoriel sur $X \simeq (\tilde{C})_a$, défini
\struck{\ill{}} \add{\uncertain{comme}} $\mathbb{W}(\mathcal{O}_X(-1))$.

L'on a \struck{\ill{}}

\begin{tikzcd}
\tilde{C} \arrow[d] & \tilde{C} \times_{C} C_a = \tilde{S} \arrow[l] \arrow[d] \\
C & C_a = S \arrow[l]
\end{tikzcd}

\begin{tikzcd}[column sep=small]
\mathrm{Pic}(C_a - a) \arrow[d, "\wr"'] & \mathrm{Pic}(C - a) \arrow[l, "\sim"'] \arrow[d, "\wr"] \\
\mathrm{Pic}(\tilde{S} - X) & \mathrm{Pic}(\tilde{C} - X) \arrow[l, "\sim"']
\end{tikzcd}

\note{à droite de ce carré, un troisième terme
$\mathrm{Pic}(\tilde{C} - X)$, avec en dessous $\mathrm{Pic}(\tilde{C})/\mathbb{Z}\eta$,
est barré d'un grand trait ; le sens des flèches horizontales et des isomorphismes
verticaux est lu sur une page chargée et reste douteux}

\page{189}
[N.B. $\mathrm{Pic}(\tilde{C} - X) \simeq \mathrm{Pic}(\tilde{C})/\mathbb{Z}.\eta$,
où $\eta$ est la classe de $X$ considéré comme diviseur dans $\tilde{C}$ ; et
comme $\mathrm{Pic}(\tilde{C}) \simeq \mathrm{Pic}(X)$ par $\tilde{C} \to X$, on
retrouve ainsi l'isomorphisme $\mathrm{Pic}(C - a) \simeq \mathrm{Pic}(X)/\mathbb{Z}.\eta$].

D'autre part on trouve aussi
\[
  \mathrm{Pic}(\hat{S} - \hat{a}) \simeq \mathrm{Pic}(\widetilde{\hat{S}} - X)
  \simeq \mathrm{Pic}(\widetilde{\hat{S}})/\mathbb{Z}\hat{\eta}
\]
où $\widetilde{\hat{S}}$ désigne l'éclaté de $\hat{S}$. Or on a un isomorphisme
évident
\[
  \widetilde{\hat{S}} \simeq \tilde{S} \times_S \hat{S} \simeq \widehat{\tilde{S}}
  \qquad \text{(complété formel de $\tilde{S}$ le long de $X$)}
\]
d'où, par le th.\ d'existence
\[
  \mathrm{Pic}(\widetilde{\hat{S}}) \simeq \mathrm{Pic}(\widehat{\tilde{S}}) ;
\]
\struck{comme par ailleurs} Ainsi, \add{l'homomorphisme}
\struck{$\mathrm{Pic}(\widetilde{\hat{S}})$}
\[
  \mathrm{Pic}(S - a) \longrightarrow \mathrm{Pic}(\hat{S} - \hat{a})
\]
est induit, en passant aux quotients par $\mathbb{Z}\eta$ resp.\ $\mathbb{Z}\hat{\eta}$,
par l'homomorphisme
\[
  \mathrm{Pic}(E) \longrightarrow \mathrm{Pic}(\hat{E})
\]

\page{190}
où $E$ est le fibré vectoriel $\tilde{C}$ sur $X$, et $\hat{E}$ son complété
formel le long de $X$, donc
\[
  \mathrm{Pic}(\hat{E}) = \varprojlim \mathrm{Pic}(E_n)
\]
où $E_n$ est le voisinage infinitésimal d'ordre $n$ de $X$ dans $E$,
\[
  E_n = (X, \mathcal{O}_E/\mathcal{J}^{n+1})
\]
$\mathcal{J}$ étant l'idéal définissant $X$ dans $E$. Notons ici
\[
  \mathcal{J}/\mathcal{J}^2 = \mathcal{O}_X(1)
\]
\[
  \mathrm{Pic}(E) \xleftarrow{\ \sim\ } \mathrm{Pic}(X)
\]
Donc $\mathrm{Pic}(S - a) \to \mathrm{Pic}(\hat{S} - \hat{a})$ est inj.\ (surj.)
ssi $\mathrm{Pic}(E) \to \mathrm{Pic}(\hat{E})$ l'est. Donc pour dire
\struck{\ill{}} \add{\ill{}} \ill{} \ill{} \ill{} (immédiatement par le
\uncertain{dernier} formule)
\marginal{(au moins dans le cas $S$ normal \ill{})}
que $\mathrm{Pic}(S - a) \to \mathrm{Pic}(\hat{S} - \hat{a})$ est injectif, i.e.\
$\mathrm{Pic}(E) \to \mathrm{Pic}(\hat{E})$ est injectif [car
$\mathrm{Pic}(X) \to \mathrm{Pic}(\hat{E})$ l'est, ayant l'inverse à gauche
$\mathrm{Pic}(\hat{E}) \to \mathrm{Pic}(E_0) = \mathrm{Pic}(X)$].
\note{tout ce passage, de « Donc » à « injectif », est serré entre les lignes et
réécrit ; l'ordre de lecture est une conjecture}

Reste : voir si $\mathrm{Pic}(E) \to \mathrm{Pic}(\hat{E})$ est bijectif. Or :
$\mathrm{Pic}(X) \simeq \mathrm{Pic}(E)$

\page{191}
\emph{il en est ainsi ssi} $\mathrm{Pic}(\hat{E}) \to \mathrm{Pic}(X)$ \emph{est
injectif}, \struck{\ill{}}
\marginal{$\mathrm{Pic}(E) \to \mathrm{Pic}(\hat{E}) \to \mathrm{Pic}(X)$,
composé $\simeq$}

Supposons par ex.\ que $\dim X = 1$ ; alors dans le système projectif
$(\mathrm{Pic}(E_n))$, les morphismes de transition sont surjectifs. D'autre
part, \add{en tous cas,} le noyau de $\mathrm{Pic}(E_n) \to \mathrm{Pic}(E_{n-1})$
($n \geq 1$) est isomorphe :
\[
  H^1(X, \mathcal{J}^n/\mathcal{J}^{n+1}) \simeq H^1(X, \mathcal{O}_X(n)) .
\]
Ainsi : \struck{Si $X$ est \ill{}}

\textbf{\emph{Proposition.}} Pour que $\mathrm{Pic}(E) \to \mathrm{Pic}(\hat{E})$
soit bijectif (ou encore $\mathrm{Pic}(S - a) \to \mathrm{Pic}(\hat{S} - \hat{a})$
bijectif) il \struck{faut} \emph{suffit} que l'on ait $H^1(X, \mathcal{O}_X(n)) = 0$
pour tout $n \geq 1$ ; cette condition est aussi nécessaire si $\dim X = 1$.

\add{Mais pour cette condition \ill{} \emph{pas vérifiée} si
$d^\circ \mathcal{O}_X(1) < g - 1$, \ill{}} Mais soit $X$
\add{\uncertain{courbe} \ill{} $\subset \mathbb{P}^2_k$, \ill{} \ill{}}
\ill{} lisse$/k$, de degré $d$, donc $g = 1 + \frac{d(d-3)}{2}$, et si
$d < g - 1$ i.e.\ $d \geq 6$, \emph{cette condition n'est jamais vérifiée}.
\marginal{Corollaire}
\note{les deux lignes « Mais pour cette condition… » et « X courbe… » sont serrées
au-dessus de la ligne suivante ; « Corollaire », encadré de deux traits
verticaux, figure en marge en face de cette fin de page}

\page{192}
\emph{Remarques.} 1) On notera que la démonstration montre que si $\dim X = 1$,
alors $\mathrm{Pic}(\hat{E}) \to \mathrm{Pic}(X)$ a un noyau qui est connexe,
\add{et \uncertain{unipotent}} \struck{\ill{}} \ill{}
$\mathrm{Pic}(S - a) \to \mathrm{Pic}(\hat{S} - a)$ induit un isom.\ sur les
groupes \emph{\uncertain{des composantes connexes}} de l'identité [qui dans ce
cas sont finis, isomorphes à $\mathbb{Z}/d\mathbb{Z}$ où $d$ est le degré
projectif de $X$] \add{et \uncertain{à un module près} \uncertain{unipotent}
\ill{}}. Le résultat semble valable si $X$ est de dim.\ quelconque, lorsque
\struck{\ill{}} la caractéristique \struck{\uncertain{résiduelle}} $p$ de $k$
est $> 0$, et que l'on raisonne \emph{\ill{}} \struck{\uncertain{projectif}}
\add{\uncertain{à pro-isomorphisme près}} ; i.e.\ le groupe de NS.\ local
\ill{} \ill{} pour $S$ et $\hat{S}$, modulo des $p$-groupes [ext.\ d'un
$p$-groupe fini par un \uncertain{groupe} \uncertain{unipotent} sur $k$] ;
utilisant le fait que les obstructions dans des $H^2(X, \underline{O}_X(n))$
finissent nécessairement par être nulles pour $n$ grand, et sont annulées par
$p$… Ceci est en accord avec la théorie du groupe fondamental …
\marginal{Regarder aussi le cas de la car.\ résiduelle nulle …}

\page{193}
2) On peut aussi regarder le hensélisé $\tilde{S}$ de $S$, et les homom.
\[
  \mathrm{Pic}(S - a) \to \mathrm{Pic}(\tilde{S} - \tilde{a}) \to \mathrm{Pic}(\hat{S} - \hat{a}) .
\]
\note{le hensélisé est noté ici d'un signe placé au-dessus du $S$, que nous
lisons comme un tilde, bien que $\tilde{S}$ désigne p.~188 l'éclaté}
Il est possible que la deuxième flèche écrite soit toujours un isomorphisme. On
le voit immédiatement, en tous cas, lorsque $X$ est de dim $1$ (\uncertain{i.e.}\
$S$ de dim $2$), par le raisonnement habituel sur l'éclaté… Dans l'exemple
précédent des courbes, il y a donc autant d'exemples où le groupe de Picard
local de $S$ et celui de son hensélisé ne sont pas isomorphes …

\section{Structure locale de $\mathbf{Pic}_{X/S}$}
\note{titre pris sur la page 194, feuille de garde de sa main : « $G^\tau$, $G^0$ etc », puis « Structure \add{loc.} de $\mathbf{Pic}_{X/S}$ », et une troisième ligne entre parenthèses, barrée, « (\ill{}, \ill{}) »}

\page{195}
\emph{Constituants de la Torsion}
\note{titre encadré, en haut à droite ; à gauche, en tête de page, $G$ au-dessus
de $S$, reliés par un trait vertical}

\uncertain{schéma} en groupes \add{\uncertain{commutatif}} \struck{\uncertain{projectif}}
sur $S$ \struck{\ill{}} noethérien \emph{local}
\struck{$G^\circ$ \uncertain{univ.}\ \uncertain{ouvert} sur $S$.}
$s \in S$ \add{pt fermé}, $s_i$ générisations immédiates.
Pour tout $t \in S$, soit $G^\circ_t$ muni de la structure \emph{abélienne}
induite, soit $T_t(G) = G_t/G^\circ_t$ et $\sigma(t) = \mathrm{rang}_{k(t)} T_t(G)$.

Implications
\begin{itemize}
  \item[(i)] \struck{$G$ \emph{plat} sur $S$,} $G^\circ$ \add{admet}
    \struck{contient} un sous-schéma \emph{abélien} $A$ ayant $G^\circ$ comme
    sous-ens.\ sous-jacent, et $G$ plat sur $S$
  \item[$\Updownarrow$]
  \item[(ii)] $G$ admet un $A$ comme dessus, et $T(G) \simeq G/A$ est \emph{plat}
    sur $S$ [N.B. il est \uncertain{diff.}\ fini sur $S$]
  \item[$\Downarrow$]
  \item[(iii)] Pour tout $t \in S$, on a \struck{rang} $\sigma(s) = \sigma(t)$
    \struck{[\ill{} $\sigma(s) \leq \sigma(s_i)$]}
  \item[$\Updownarrow$]
  \item[(iv)] Pour tout $i$, \ill{}
    $\sigma(s) = \sigma(s_i)$ [on sait $\sigma(s) \leq \sigma(s_i)$]
\end{itemize}
\note{la flèche entre (iii) et (iv) est double, sa pointe supérieure
surchargée ; la phrase suivante la conditionne}
\add{\ill{} i.e.\ les fibres de $G$ sont toutes de même dim.}
Si $G^\circ$ univ.\ ouvert sur $S$, on a (iii) $\Leftrightarrow$ (iv).
Si \add{de plus} $S$ est \emph{réduit}, \emph{alors} les conditions envisagées
sont équivalentes.

On sait que (théorie des groupes \uncertain{quotients} \add{plats}) que si $A$
existe, alors $G \to G/A$ est fid.\ plat et projectif, et que alors $G$ plat sur
$S$ ssi $G/A$ plat sur $S$, ce qui démontre l'équivalence de (i) et (ii). De
plus, si $A$ existe, alors on a pour tout $t \in S$,
\[
  T_t(G) \simeq T(G)_t \qquad \text{où } T(G) = G/A.
\]
Donc (ii) $\Rightarrow$ (iv), a fortiori (ii) $\Rightarrow$ (iii).

Pour (iii) $\Leftrightarrow$ (iv), on \struck{\ill{}} \add{\uncertain{utilise le}}
(Supposons $G^\circ$ univ.\ ouvert sur $S$, i.e.\ les fibres de $G$ sont toutes
de même dim.)

\textbf{\emph{Lemme}} (La fonction $\sigma$ est \emph{semi-continue}
\emph{supérieurement}. Comme elle est évidemment constructible (cf l'argument
précédent …] on est ramené à prouver, dans le cas où $S$ est le spectre
d'un anneau de val.\ discrète, s'il pt fermé, \ill{} que
$\sigma(s) \geq \sigma(s')$. Soit $G^\circ$ muni de la structure \uncertain{induite}
de celle de $G^\circ_{s'}$, donc $G^\circ$ est un \emph{sous-groupe} de $G$
\emph{plat}$/S$, et on peut prendre le groupe quotient $Q = G/G_0$, qui est
fini sur $S$. Alors
\[
  \sigma(s) \geq \mathrm{rang}_{\mathcal{O}}[Q_s : k(s)] \geq
  \mathrm{rang}_{\mathcal{O}}[Q_{s'} : k(s')] = \sigma(s')
\]
d'où l'inégalité voulue. Notons que la première inégalité est une égalité ssi
$G^\circ_s$ est réduit, i.e.\ $G^\circ$ est \emph{abélien} sur $S$, et que la
deuxième l'est ssi $Q$ est plat sur $S$. / Cela prouve \struck{\ill{}}
\add{\uncertain{aussi}} que dans le cas $S = \mathrm{Spec}(V)$, on a
(iv) $\Rightarrow$ (ii). C'est aussi vrai dès que $S$ est seulement réduit, En
effet, on \uncertain{sait} \ill{}
\marginal{\ill{} trouve facilement un \uncertain{sous}-schéma en
\uncertain{groupes} \ill{} \ill{} schéma $S' \to S$ \uncertain{radiciel} \ill{}
le pt \ill{} d'un $A$. Alors $S' \to S$ est un \ill{} \ill{}, i.e.\
\emph{propre}. Le critère \uncertain{valuatif} \ill{} \ill{} les raisonnements
précédents s'appliquent \ill{} $T = G/A$. Et \ill{} $T \to S$ \ill{} \ill{}
schéma de Hilbert de $G/S$ ; \ill{} universel correspondant \ill{} \ill{}
\ill{}, et surjectif, \ill{} \ill{} \ill{} \uncertain{critère} \ill{}
$S = \mathrm{Spec}(V)$, \ill{}. D'où résultat, l'existence de $A$, \ill{} en
vertu de (iv), [\ill{} critère val.\ \ill{}]}
\note{la note marginale, en deux colonnes écrites de bas en haut le long du bord
gauche, poursuit la dernière phrase de la page ; elle n'est lue que par
fragments}

\page{197}
\note{feuillet au crayon ; le texte au-dessus de l'accolade se rattache aux
conditions (iii), (iv) de la p.~195}
\textbf{\emph{Corollaire}} (\struck{\ill{}} \uncertain{fibres} de $G$
\uncertain{connexes} et \uncertain{de même} dimension). Les conditions (iii) et
(iv) \uncertain{équivalent} \ill{} à la suivante, \uncertain{où} $n$ est un
entier fixé \struck{\ill{}} multiple de $\sigma(t)^2$, et
${}_nG = \mathrm{Ker}(G \xrightarrow{n} G)$ :
\begin{itemize}
  \item[(v)] $\mathrm{rang}_{k(s)}\, {}_nG_s = \mathrm{rang}_{k(s_i)}\, {}_nG_{s_i}$
    pour tout $i$,
\end{itemize}
ou
\begin{itemize}
  \item[(v bis)] $\mathrm{rang}_{k(s)}\, {}_nG_s = \mathrm{rang}_{k(t)}\, {}_nG_t$
    pour tout $t \in S$.
\end{itemize}

\note{en marge droite, deux fragments isolés : « $H^1(W_{nn}$ » et
« $F^n V^n = 0$ »}

Dans un cadre :
\[
  \mathrm{long}\, H^1(\bar{X}_s, \mathbb{Z}/p^n)
  + \mathrm{long}\, H^1(\bar{X}_s, \mu_{p^n})
\]
\[
  + \;\struck{\ill{}}\ \mathrm{long}_{W_n^\sigma[F,V]}\, H^1(\bar{X}_s, W_{n,n})
\]
partie \uncertain{multiplicative} \quad partie additive \add{\uncertain{unipotente}}

\note{sous le cadre, trois signes $\Vert$ renvoient chaque terme à une
interprétation :}
\[
  \mathrm{Hom}(\pi_1(\bar{X}_s), \mathbb{Z}/p^n) , \qquad
  {}_{p^n}\mathrm{Pic}(\bar{X}_s) ,
\]
\[
  \mathrm{Ker}\bigl(H^1(\bar{X}_s, W_n) \xrightarrow{F^n} H^1(\bar{X}_s, W_n)\bigr)
\]

\section{Sous-groupe $G^\tau$}
\note{titre encadré en tête de la p.~200}

\page{200}
Soit $G$ un schéma en groupes sur le \emph{corps} $k$, localement de type fini sur
$k$. Alors :

\struck{1°) La composante connexe de l'élément neutre $e$ est de type fini sur $k$,
et géométriquement irréductible}
\note{cet énoncé est barré d'un trait ondulé et remplacé par 1°) et 2°)}

1°) Les composantes connexes \add{$G_i$} de $G$ sont de type fini sur $k$, et
irréductibles

2°) La composante connexe \add{\struck{\ill{}} $G_0$} de $e$ est géom.\ irréductible,
et c'est un sous-groupe invariant \struck{\ill{}}

3°) \add{Les $G_i$ sont stables \uncertain{par} $G_0$.} \struck{Si $k$ est alg.\
clos, alors les \struck{\ill{}} composantes connexes de $G$ sont les
\uncertain{translatés} de $G^\circ_0$, où $a \in G(k)$}
\add{Si} $G_i$ contient un pt rat.\ sur $k$, alors $G_i = G_0 a$
\note{le passage barré de 3°) est entouré et rayé de traits ondulés ; « Si $G_i$
contient… » est écrit entre ses lignes}

4°) Le \struck{groupe} \add{schéma quotient} $G/G^\circ_0$ existe et est discret
i.e.\ localement fini sur $k$, $G \to G/G^\circ_0$ est un morphisme
\struck{\ill{}} fid.\ plat. $G/G_0$ est un $k$-groupe, séparé et loc.\ fini.
\struck{\ill{}}

\struck{La composante connexe $G^\circ_{\bar{k}}$ de $e_{\bar{k}}$ de $G_{\bar{k}}$ \ill{}}

Supposons d'abord $k$ alg.\ clos. \struck{Soit $G_0$} \struck{définissons} Comme
$G_0$ est connexe, il en est de même de $G_0 \times G_0$ (référence …), donc
$G_0 \times$ \struck{$G_0$} $G_0 \to G$ induit $G_0 \times G_0 \to G_0$, i.e.\
$G_0$ est un sous-groupe. Soit $U$ un \add{\uncertain{ouvert}} voisinage
\ill{} de $e$, je dis que
\[
  U^2 = G_0
\]
\marginal{$\times$ \uncertain{utilisant}, $G_0$ \ill{} \uncertain{irréductible}
\ill{} de type fini}
\uncertain{Et} \ill{} $G_0$ est quasi-compact …
\note{l'argument se poursuit au-delà de ce lot}

\end{document}
